\chapter{流体静力学}

\begin{yijubox}
E1 \texttt{8-1考情分析.pptx} slide33：第二章流体静力学＝计算题基础入门、重点掌握方法。\\
E1 slide25：真题重复率很高，围绕参考书复习即可，总体难度中等。\\
E2 \texttt{9-1考点分析.pptx}：计算题占 50\% 分值；第 1--3 章偏概念，掌握基本公式与解释即可。\\
E3 官方书＝赵琴《流体力学与流体机械》，初试只考前 8 章；真题大部分源于课后习题改编，
第二章课后习题（习题 2.1--2.17）务必逐题过关。\\
\textbf{本章结论}：静压强特性、平衡微分方程、基本方程与压强表示是"概念＋小计算"，
必须在选择、填空里拿满分；平面／曲面总压力（含压力体）是本章唯一能出计算大题的地方，
也是第 4 章动量方程、第 5 章水击的前置工具——本章方法不熟，后面整条计算链都断。
\end{yijubox}

\section{流体静压强特性}

\begin{kaodian}{4}{绝对静止与相对静止}
\textbf{是什么：}\textbf{绝对静止}＝流体整体对地球没有相对运动，此时流体所受的质量力
\textbf{只有重力}；\textbf{相对静止}（相对平衡）＝流体整体对地球有相对运动，但
\emph{流体质点之间没有相对运动}，各质点与容器的相对位置也不改变。

\textbf{两者的共同点（考点所在）：}两种状态下流体质点之间都\emph{没有相对运动}，
黏性都体现不出来，流体内部都\emph{不存在切应力}，表面力都只剩压力。
所以流体静力学的基本概念、欧拉平衡微分方程和等压面理论对两种状态\emph{同样适用}；
区别只有一处——所受质量力不同：绝对静止只有重力，相对静止还有惯性力。

\textbf{常考题型：}名词解释／填空（"相对静止是指\underline{\hspace{3cm}}"）；
选择（"在相对平衡的液体中，黏性表现为（\quad）"$\to$ 不表现）；
简答（"绝对静止与相对静止有何异同"）。

\textbf{重要度 \xstarfull{4}}\quad\yiju{E1 slide33「第二章＝计算题基础入门、重点掌握方法」；E1 slide16 填空考名词定义；相对平衡题需先判质量力，见 2.4 节}\quad\nandu{易}

\begin{banfa}
判断题万用一句话：\textbf{"静止流体不能承受切应力，只能承受压力"}。
凡选项出现"静止流体内部有切应力""静止流体黏性起作用"$\to$ 一律错。
判断"是绝对静止还是相对静止"，只看\emph{对地球有无整体运动}、\emph{质量力里有没有惯性力}，
不要看容器里液体表面是平还是斜。
\end{banfa}
\end{kaodian}

\begin{kaodian}{4}{流体静压力（静压强）及其两大特性与单位}
\textbf{是什么：}在静止流体中不存在切应力，单位面积上所受到的垂直于该表面的力称为
\textbf{流体静压力}（简称静压力、静压强）。

\begin{gongshi}{静压强的定义式（式 2.1、2.2）}
\[ p=\lim_{\Delta A\to0}\frac{\Delta F}{\Delta A} \]
\textbf{含义：}取包含某点的微元面积 $\Delta A$，其上的法向压力为 $\Delta F$，
令面积趋近于零所得极限即为该点的\textbf{点静压强}，单位 $\mathrm{Pa}=\mathrm{N/m^2}$。\\
\textbf{适用条件：}连续介质假设成立；静止（或相对静止）流体；$\Delta F$ 必须是
\emph{法向}压力，不含任何切向分量。
\end{gongshi}

\textbf{两大特性（必背，历年反复考）：}
\begin{center}
\begin{tabularx}{\textwidth}{@{}lXX@{}}
\toprule
特性 & 表述 & 物理解释 \\
\midrule
特性一（方向） & 静压强方向必然沿作用面的\textbf{内法线方向}，即\emph{垂直并指向}作用面 & 静止流体无切应力，表面力只剩法向压力；流体几乎不能承受拉力，故只能是压力 \\
特性二（大小） & 静止流体中\textbf{任一点}处的压强大小\emph{与其作用面方位无关}，同一点各方向的静压强大小相等 & 取微小四面体 $MABC$ 对 $x,y,z$ 三轴分别列力平衡，得 $p_x=p_y=p_z=p_n$（式 2.1、2.2） \\
\bottomrule
\end{tabularx}
\end{center}

由特性二可写出压强只是位置的函数：
\[ p=p(x,y,z) \]
（式 2.3）——正因为如此，才能说"静止流体中\textbf{某一点}的压强"，
也才能在 2.2 节把压强当成坐标的连续可微函数来求偏导数。

\textbf{单位与换算（填空常考）：}国际单位制 $\mathrm{Pa}$（$1\ \mathrm{Pa}=1\ \mathrm{N/m^2}$），
还有 $\mathrm{kPa}$、$\mathrm{MPa}$、$\mathrm{bar}$（$1\ \mathrm{bar}=10^5\ \mathrm{Pa}$）；
工程上常用\textbf{液柱高度}作压强单位：米水柱（$\mathrm{mH_2O}$）、毫米汞柱（$\mathrm{mmHg}$），
并记 $1\ \mathrm{atm}=760\ \mathrm{mmHg}$。

\textbf{常考题型：}选择（"静压强方向必沿（外法线／内法线）"；
"同一点各方向静压强（相等／不等）"）；填空（静压强的单位、四大特性辨析）。

\textbf{重要度 \xstarfull{4}}\quad\yiju{E1 slide33「第二章＝计算题基础入门、重点掌握方法」；E1 slide16「选择/判断考概念辨析、填空考名词定义与小计算」；E2「第 1--3 章偏概念，掌握基本公式与解释即可」}\quad\nandu{易}

\begin{yicuo}
\textbf{"内法线"是唯一正确答案。}作用面的法线有两条（外法线指向流体外部、内法线指向
流体内部），静压强指向受作用的\emph{流体或壁面}内部，所以是内法线。
把特性一说成"沿作用面的法线方向"（不区分内外）在简答题里要扣分。\\
另一处：特性二说的是"\emph{同一点}上不同方向的压强相等"，不是"不同点压强相等"——
不同点压强一般\emph{不}相等，否则 2.3 节的 $p=p_0+\rho g h$ 就不成立了。
\end{yicuo}
\end{kaodian}

\section{流体平衡微分方程}

\begin{kaodian}{5}{欧拉平衡微分方程与全微分形式}
\textbf{是什么：}在静止流体中取微元直角六面体（中心 $M$，边长 $\dd x,\dd y,\dd z$），
分别列出三个方向的力平衡，即得欧拉平衡微分方程。

\begin{gongshi}{欧拉平衡微分方程（式 2.4、2.5）}
\[
\begin{cases}
\dfrac{\partial p}{\partial x}=\rho f_x\\[2mm]
\dfrac{\partial p}{\partial y}=\rho f_y\\[2mm]
\dfrac{\partial p}{\partial z}=\rho f_z
\end{cases}
\qquad\text{矢量式：}\quad \grad p=\rho \bF
\]
（式中 $f_x,f_y,f_z$ 为单位质量力 $\bF$ 的三个分量，$\grad$ 为哈密顿算子。）
\textbf{含义：}表面力（压力）与质量力在三个坐标方向\emph{逐一平衡}；
写全了是 $f_x-\dfrac{1}{\rho}\dfrac{\partial p}{\partial x}=0$ 等三式，
移项即得上式（两种写法等价，考试都算对）。\\
\textbf{适用条件：}\textbf{任何静止流体}——\emph{不分}绝对静止还是相对静止，
\emph{不分}可压缩还是不可压缩，\emph{不分}均质还是非均质，是本章最普适的方程。
欧拉于 1775 年提出，是流体力学第一个微分方程。
\end{gongshi}

\begin{gongshi}{全微分形式与力的势函数（式 2.6）}
\[ \dd p=\rho\,(f_x\,\dd x+f_y\,\dd y+f_z\,\dd z) \]
若质量力有势，即 $f_x=-\dfrac{\partial W}{\partial x}$、$f_y=-\dfrac{\partial W}{\partial y}$、
$f_z=-\dfrac{\partial W}{\partial z}$，则 $\dd p=-\rho\,\dd W$，积分得
\[ p=-\rho W+C \]
\textbf{含义：}把三式分别乘 $\dd x,\dd y,\dd z$ 再相加，左边恰是压强的全微分 $\dd p$。\\
\textbf{适用条件：}流体平衡；$\rho$ 在积分路径上为常数（不可压缩均质流体）时才能直接积分；
有势质量力是能积分成 $p=-\rho W+C$ 的\emph{充分条件}。
\end{gongshi}

\textbf{常考题型：}选择（"重力场中静压强微分方程"$\to\dfrac{\partial p}{\partial z}=-\rho g$，
即只保留 $z$ 方向一项）；填空（写出全微分形式 $\dd p=\rho(f_x\dd x+f_y\dd y+f_z\dd z)$）；
简答（欧拉平衡微分方程的物理意义＝单位质量流体的质量力与压力合力相平衡）。

\textbf{重要度 \xstarfull{4}}\quad\yiju{E4 2022 年 818 单选第 5 题（重力场静压强微分方程 $\dd p$）；E5 题库选择第 3 题（体积力有势）；E1 slide33「第二章＝计算题基础入门、重点掌握方法」}\quad\nandu{中}

\begin{banfa}
\textbf{①记方程的"三步来源"}：取六面体 $\to$ 六个面上用泰勒展开写出压强 $\to$
表面力之和 $=$ 质量力，即得式（2.4）。考简答时写这三步就是标准答案。\\
\textbf{②判断质量力是否"有势"}的速判法：看这个力做的功是否与路径无关。
重力、离心惯性力、等加速直线运动的惯性力\emph{都有势}；
所以本章的相对平衡问题都能用 $\dd p=\rho(f_x\dd x+f_y\dd y+f_z\dd z)$ 积分求解。\\
\textbf{③选坐标系先写 $f_x,f_y,f_z$}：例如容器沿水平方向以加速度 $a$ 运动，
取 $x$ 轴沿运动方向、$z$ 轴铅垂向上，则 $f_x=-a$、$f_y=0$、$f_z=-g$，
代进方程立刻能积出压强分布。\emph{这一步是相对平衡题的第一分点}。
\end{banfa}
\end{kaodian}

\begin{kaodian}{4}{等压面及其性质}
\textbf{是什么：}\textbf{压强相等的点组成的面}称为等压面，可以是平面也可以是曲面。
液体的自由表面（与气体接触的面）就是特殊的等压面——其上各点压强都等于
气体压强 $p_0$。

\begin{gongshi}{等压面方程与性质（式 2.7、2.8）}
\[ f_x\,\dd x+f_y\,\dd y+f_z\,\dd z=0 \qquad\text{矢量式：}\quad \bF\cdot\dd\bq=0 \]
\textbf{含义：}在等压面上 $\dd p=0$，代入全微分形式即得；矢量式表明
\textbf{质量力与等压面互相垂直}——这是等压面最重要的性质。\\
\textbf{适用条件：}同一种连续流体，且等压面\emph{不能穿越密度不同的两种流体的分界面}。
\end{gongshi}

\textbf{三条推论（判断、选择题几乎年年用）：}
\begin{center}
\begin{tabularx}{\textwidth}{@{}lXX@{}}
\toprule
质量力情况 & 等压面形状 & 应用 \\
\midrule
只有重力（绝对静止） & 质量力铅垂向下，等压面为\textbf{一簇水平面} & 静水中同一水平面上压强相等 \\[1mm]
重力 $+$ 水平惯性力（等加速直线运动） & 等压面为\textbf{一簇倾斜平面}，斜率由 $a/g$ 决定 & 油罐车、洒水车自由面倾斜 \\[1mm]
重力 $+$ 离心惯性力（等角速度旋转） & 等压面为\textbf{一簇旋转抛物面} & 离心机、离心铸造、离心泵入口 \\
\bottomrule
\end{tabularx}
\end{center}

\textbf{等压面的选取法（测压管、压差计题的命门）：}在同一种\emph{静止连通的同种液体}中
找同一高度、且\emph{两部分液体相同}的位置，那里压强相等。
实用口诀：\textbf{从最低的同一液体分界面倒着往两边推}。

\textbf{常考题型：}选择（"等压面与质量力的关系"$\to$ 垂直）；
填空（写出等压面方程）；简答（等压面的三个成立条件：同种流体、静止连通、水平面取同高度）。

\textbf{重要度 \xstarfull{4}}\quad\yiju{E1 slide33「第二章＝计算题基础入门、重点掌握方法」；E1 slide16 选择/判断考概念辨析；等压面是 2.3.3 测压管、2.4 相对平衡两类题的共同入口（赵琴书 2.2.3 节）}\quad\nandu{易}

\begin{yicuo}
\begin{enumerate}
\item "等压面就是水平面"——\textbf{错}，只有在质量力只有重力时才是。
旋转容器里的等压面是抛物面，加速容器里是斜面。
\item "等压面必须是平面"——\textbf{错}，也可以是曲面。
\item 在两种不同液体的\emph{分界面}上取"等压面"——如果分成同一高度却不是同种液体，
  两侧压强\emph{不}相等，绝对不能取。
\end{enumerate}
\end{yicuo}
\end{kaodian}

\section{重力场中流体静压强分布}

\begin{kaodian}{5}{流体静力学基本方程与测压管水头}
\textbf{是什么：}质量力只有重力、不可压缩均质流体时，欧拉平衡微分方程的积分结果。

\begin{gongshi}{静压强基本方程（式 2.9--2.11）}
\[ z+\frac{p}{\rho g}=C \qquad\qquad
   z_1+\frac{p_1}{\rho g}=z_2+\frac{p_2}{\rho g} \qquad\qquad
   p=p_0+\rho g h \]
（$z$ 由某水平基准面铅垂向上量取，$h$ 为该点在液面以下的淹深。）
\textbf{含义：}由 $f_x=f_y=0$、$f_z=-g$ 得 $\dfrac{\dd p}{\dd z}=-\rho g$，
积分即得 $z+\dfrac{p}{\rho g}=$ 常数；再把一点取在液面（压强 $p_0$）即得
$p=p_0+\rho g h$——\emph{液面传递的压强与液柱自重压强之和}。\\
\textbf{适用条件：}\textbf{重力}作用下、\textbf{同一种}（均质）\textbf{不可压缩}
（$\rho$ 常数）\textbf{静止}流体，且\emph{互相连通}；不同密度分层时要逐层用，
每换一种液体就要重新积分一次。
\end{gongshi}

\textbf{两项意义（简答必背）：}
\begin{center}
\begin{tabularx}{\textwidth}{@{}lXX@{}}
\toprule
 & 物理意义（能量） & 几何意义（水头） \\
\midrule
$z$ & 单位重量流体对基准面的\textbf{位能}（比位能） & \textbf{位置水头} \\[1mm]
$\dfrac{p}{\rho g}$ & 单位重量流体具有的\textbf{压能}（比压能） & \textbf{压强水头}（压力水头） \\[1mm]
$z+\dfrac{p}{\rho g}$ & 单位重量流体的\textbf{总势能}（比势能） & \textbf{测压管水头}（静水头） \\
\bottomrule
\end{tabularx}
\end{center}

由此得本章最常考的两条结论：
\begin{itemize}
\item 静止流体中\textbf{各点总势能相等}（能量守恒在静止流体中的表现）；
\item 静止流体中\textbf{各点测压管水头相等}，所以\textbf{测压管水头线是一条水平线}。
\end{itemize}

$p=p_0+\rho g h$ 的三条推论（赵琴书口径）：
① 重力作用下均质流体内部静压强与深度 $h$ 呈\textbf{线性}关系
（所以水坝设计成上窄下宽）；
② 任一点静压强由"液面压强 $p_0$"和"液柱压强 $\rho g h$"两部分组成，
\textbf{深度相同处压强相等}；
③ 边界（液面）上压强的变化将均匀传到流体每一点，即\textbf{帕斯卡定律}
（这就是水压机、液压传动的原理）。

\textbf{常考题型：}选择（"测压管水头线是（水平／倾斜）线"；"两点等深，压强大小关系"）；
填空（位置水头、压强水头、测压管水头的定义）；计算（复式测压计读数换算，
如赵琴书习题 2.3、习题 2.5 三种互不相混液体求三根测压管液面高度）。

\textbf{重要度 \xstarfull{5}}\quad\yiju{E1 slide33「第二章＝计算题基础入门、重点掌握方法」；E4 2022 年 818 单选第 3 题考测压管水头线的沿程变化；E1 slide16 填空考名词定义与小计算；E2「第 1--3 章掌握基本公式与解释即可」}\quad\nandu{中}

\begin{banfa}
\textbf{①"水平面等压"只用一次、用对地方}：只有\emph{连通}的\emph{同种}液体中
同一水平面才等压；一旦穿过两种液体的分界面，就必须改用在分界面处压强连续这个条件。\\
\textbf{②复式测压计（多段流体）三步走：}(a) 从已知压强那一端出发；
(b) 每向下走 $h$ 就加 $\rho g h$、向上走 $h$ 就减 $\rho g h$（\emph{谁在走就用谁的密度}）；
(c) 落点在自由液面或大气则令其等于 $p_a$，解方程。这是赵琴书例 2.1、习题 2.2--2.4 的统一解法。\\
\textbf{③$p=p_0+\rho g h$ 的 $h$ 一定要"从自由液面量起"}，
不能从容器底面量，也不能把两种液体混用一个 $h$。
\end{banfa}

\begin{yicuo}
\begin{enumerate}
\item $p=\rho g h$ 中的 $h$ 是\textbf{淹深}（自由液面向下量），不是坐标 $z$，
  也不是容器高度；公式里的 $p$ 若写成 $p_0+\rho g h$ 则默认 $p_0$ 是液面压强。\\
\item "深度相同处压强相等"的前提是\textbf{同一种连通液体}。
  在分层的油水容器里，同深度不同侧压强并不一定相等。\\
\item 测压管水头 $z+\dfrac{p}{\rho g}$ 与\textbf{总水头} $z+\dfrac{p}{\rho g}+\dfrac{v^2}{2g}$
  是两个东西：静止流体 $v=0$ 时二者才相等。第 4、5 章常把二者混在一起设陷阱。
\end{enumerate}
\end{yicuo}
\end{kaodian}

\begin{kaodian}{4}{压强表示法：绝对压强、表压强与真空度}
\textbf{是什么：}同一压强可有不同的基准（零点）。

\begin{gongshi}{三种压强的换算（式 2.12、2.13）}
\[ p_{\text{abs}}\ \text{（绝对压强）},\qquad
   p_M=p_{\text{abs}}-p_a\ \text{（表压强）},\qquad
   p_v=p_a-p_{\text{abs}}\ \text{（真空度）} \]
\textbf{含义：}以\emph{绝对真空}为零点的称\textbf{绝对压强} $p_{\text{abs}}$，
它反映压强的物理本质；以\emph{当地大气压}为零点的称\textbf{相对压强}。
相对压强为正时称\textbf{表压强} $p_M$（恒正）；为负时用\textbf{真空度} $p_v$ 表示，
使读数取正，$p_v$ 越大说明真空状态越显著。\\
\textbf{适用条件：}$p_a$ 必须取\emph{当地}大气压（随地区、季节、气候变化）；
工程问题中不特别说明时压强均指\textbf{相对压强}。
\end{gongshi}

三条换算关系要背成一句话：
\[ p_M=p_{\text{abs}}-p_a,\qquad p_v=p_a-p_{\text{abs}}=-p_M,\qquad
   p_{\text{abs}}=p_a+p_M=p_a-p_v \]
坐标系排列：\textbf{绝对真空（0）$\to$ 真空区（$p_v>0$）$\to$ 当地大气压 $p_a$ $\to$
表压区（$p_M>0$）}，$p_a$ 是分界。

\textbf{常考题型：}选择（"真空度越大说明绝对压强越（小／大）"$\to$ 越小）；
填空（$p_a=98\ \mathrm{kPa}$ 时某点 $p_{\text{abs}}=30\ 772\ \mathrm{Pa}$，则真空度为 $67\ 228\ \mathrm{Pa}$）；
计算（赵琴书例 2.1：由水银压力计读数求水管 $A$ 的真空度与绝对压强）。

\textbf{重要度 \xstarfull{4}}\quad\yiju{E1 slide33「第二章＝计算题基础入门、重点掌握方法」；E4 2022 年 818 单选第 7 题考虹吸管最高处压强与大气压的关系（判断压强低于／高于大气压需用本节的三种压强转换）；E1 slide16 填空考小计算}\quad\nandu{易}

\begin{banfa}
\textbf{①一句话判真空：}绝对压强小于当地大气压就是真空，真空度 $=p_a-p_{\text{abs}}>0$。
题目问"真空度"，答数必须是\textbf{正数}；若算出负数，说明该点其实是表压，要换公式。\\
\textbf{②虹吸管最高处的经典结论：}$p_{\text{abs}}<p_a$，即管内出现\textbf{真空}，
所以选项里说"最高处压强大于大气压"必错。\emph{这一分靠本节概念就能拿到。}\\
\textbf{③单位互换速记：}$1\ \mathrm{atm}=760\ \mathrm{mmHg}=10.33\ \mathrm{mH_2O}=101.3\ \mathrm{kPa}$；
水的压强水头每 $1\ \mathrm{m}$ 对应 $9.8\ \mathrm{kPa}$，水银每 $1\ \mathrm{mm}$ 对应 $133.3\ \mathrm{Pa}$。
\end{banfa}
\end{kaodian}

\begin{kaodian}{4}{静压强的测量：测压管、U 形测压管与压差计}
\textbf{是什么：}常用测压方式有弹簧金属式、电测式和\textbf{液位式}三类；
液位式的原理就是静压强基本方程，主要设备是测压管、U 形测压管和压差计。

\begin{gongshi}{测量设备与读数公式}
\[
\text{测压管：}\ p_B=\rho g h,\quad p_{B,\text{abs}}=p_a+\rho g h;\qquad
\text{U 形测压管（气体）：}\ p_0=\rho_{\text{Hg}}g\,\Delta h
\]
\textbf{含义：}\textbf{测压管}是一端接被测点、另一端竖直通大气的连通管，
管内液柱高度即被测点的\textbf{相对压强}；\textbf{U 形测压管}内装水银等工作流体，
靠读出两侧液面高差 $\Delta h$ 反算压强。\\
\textbf{适用条件：}测压管宜测不太大的相对压强（相对压强 $0.1$ 个大气压时水柱已高 $1\ \mathrm{m}$，
再大不便测读），且\textbf{管径 $d\ge5\ \mathrm{mm}$} 以免表面张力影响读数；
测气体时用 U 形管，测压管不适于测气体压强。
\end{gongshi}

\textbf{压差计：}U 形管水银压差计测两点压强差，在等压面 $D$--$D$ 上有
\[ p_1+\gamma_{\text{油}}h_1=p_2+\gamma_{\text{油}}h_2+\gamma_{\text{Hg}}\Delta h \]
整理得 $\Delta p=p_1-p_2=\gamma_{\text{油}}(h_2-h_1)+\gamma_{\text{Hg}}\Delta h$
（注意 $\gamma=\rho g$）。若两侧\emph{同种}液体，则化简为 $\Delta p=(\gamma_{\text{Hg}}-\gamma)h$。

\textbf{常考题型：}填空（"测压管直径一般不小于 $5\ \mathrm{mm}$"；U 形压差计 $\Delta p$ 表达式）；
计算（赵琴书例 2.1 水银压力计求真空度与绝对压强；习题 2.2 上部有空气的压差计求两点压强差；
习题 2.5 三种不相混液体求各测压管液面高度）。

\textbf{重要度 \xstarfull{4}}\quad\yiju{E4 2015 年试题填空考 U 形水银压差计；E1 slide16「填空考名词定义与小计算」；E3「真题大部分源于课后习题改编」——赵琴书习题 2.1--2.5 全属此类}\quad\nandu{中}

\begin{banfa}
\textbf{统一解法（管它几段流体、几根管子）：}
① 选一个\emph{两侧同种液体的最低位置}作等压面；
② 从等压面分别向两侧"走到"已知压强点，\emph{向下加 $\rho g h$、向上减 $\rho g h$}；
③ 令两侧相等，解未知量。\\
\textbf{易错防护：}走的过程中每换一种液体就必须换一次 $\rho$；
U 形管上部是\textbf{空气}时，空气柱自重 $\rho_{\text{air}}gh$ 可以忽略，直接令两侧水柱等压。
\end{banfa}
\end{kaodian}

\section{流体的相对平衡}

\begin{kaodian}{3}{等加速直线运动容器中的相对平衡}
\textbf{是什么：}容器以等加速度 $\bm{a}$ 沿水平方向运动，液体随容器一起运动、
质点之间无相对运动。液体除受重力外还受\emph{水平惯性力}，此时用欧拉平衡微分方程分析。

\begin{gongshi}{等加速直线运动容器的自由面与压强分布}
\[
\tan\theta=\frac{a}{g},\qquad z_s=-\frac{a}{g}x+C,\qquad
p=p_0+\rho g\left(\frac{a}{g}x_0+h\right)
\]
\textbf{含义：}取 $x$ 轴沿运动方向、$z$ 轴铅垂向上，则 $f_x=-a$、$f_z=-g$，
代入 $\dd p=\rho(f_x\dd x+f_z\dd z)$ 积分得等压面为\textbf{倾斜平面}，
斜率由 $\tan\theta=a/g$ 决定；自由面随加速度增大而倾斜更陡。\\
\textbf{适用条件：}容器做\textbf{等加速直线运动}、液体已达相对静止（不晃动）、
不可压缩均质液体；$x$ 轴必须取在运动方向上。
\end{gongshi}

\textbf{关键结论：}求出自由面方程后，\textbf{液体内任一点的压强仍可用 $p=p_0+\rho g h$
计算}，只是 $h$ 是"该点到\emph{倾斜}自由面的铅垂距离"。

\textbf{常考题型：}选择（自由面斜率与 $a/g$ 的关系；$a$ 增大时倾角变化）；
计算（给定加速度与容器尺寸，求自由面方程、液面倾斜后的最高／最低点、某点压强）。

\textbf{重要度 \xstarfull{3}}\quad\yiju{E1 slide33「第二章＝计算题基础入门、重点掌握方法」；E4 历年真题中未见直接题号，仅章节级依据；E2「第 1--3 章掌握基本公式与解释即可」}\quad\nandu{中}

\begin{banfa}
\textbf{三步走：}① 先建系、写 $f_x=-a$、$f_z=-g$；② 用 $\dd p=0$ 求等压面／自由面方程，
拿到 $\tan\theta=a/g$；③ 用"\emph{体积守恒}定自由面位置"（液体总量不变，
倾斜后一侧上升的体积等于另一侧下降的体积），再代 $p=p_0+\rho g h$ 求某点压强。
第③步是这类题最容易漏的得分点。
\end{banfa}
\end{kaodian}

\begin{kaodian}{3}{等角速度旋转容器中的相对平衡}
\textbf{是什么：}半径为 $R$ 的圆筒容器绕铅垂轴以恒定角速度 $\omega$ 旋转，
液体在器壁带动下最终以同一角速度旋转，自由面由水平面变成\textbf{旋转抛物面}。
质点受重力和\emph{离心惯性力}。

\begin{gongshi}{旋转容器内的压强、等压面与自由面方程（式 2.14--2.20）}
\[
f_x=\omega^2x,\quad f_y=\omega^2y,\quad f_z=-g
\]
\[ p=p_0+\rho g\left(\frac{\omega^2r^2}{2g}-z\right)\ \text{（式 2.15）} \]
\[
z=\frac{\omega^2r^2}{2g}+\frac{p-p_0}{\rho g}\ \text{（等压面方程）},\qquad
z_0=\frac{\omega^2r^2}{2g}\ \text{（自由面方程）}
\]
\[
\Delta H=\frac{\omega^2R^2}{2g}\ \text{（最大超高）},\qquad
V=\pi R^2\frac{\Delta H}{2}\ \text{（回转抛物体体积）}
\]
（坐标原点取在圆筒轴心线与\emph{原}液面交点上，$z$ 轴与轴心线重合向上，$r=\sqrt{x^2+y^2}$。）
\textbf{含义：}离心惯性力随半径线性增大，故等压面是\emph{抛物面}；
液面最大超高 $\Delta H$ 与 $\omega^2$、$R^2$ 成正比；
抛物面体内的体积恰为"高为 $\Delta H$ 的圆柱体体积的一半"。
\textbf{压强公式（式 2.20）：}$p=p_0+\rho g(z_0-z)=p_0+\rho g h$——
仍然可用典型静压强公式，$h$ 为点到\emph{自由面}的淹深。\\
\textbf{适用条件：}液体已达等角速度旋转（相对静止）、不可压缩均质液体、
原点取在自由面中心；若原点改取容器底面与转轴交点，
积分常数 $C$ 必须重新用边界条件确定（赵琴书例 2.2）。
\end{gongshi}

\textbf{两个必须记住的"水平面压强分布"差异：}
\begin{center}
\begin{tabularx}{\textwidth}{@{}lXX@{}}
\toprule
状态 & 同一水平面内压强 & 工程应用 \\
\midrule
绝对静止 & 各点压强\textbf{相等} & 常规测压、水坝 \\
绕铅直轴等角速度旋转 & 压强随半径 $r$ 变化，\textbf{轴心最低、边缘最高} & 离心铸造机 \\
\bottomrule
\end{tabularx}
\end{center}

\textbf{开口位置决定压强分布（易混）：}上盖\textbf{中心}开小孔通大气时，
上盖内表面 $p=\frac{1}{2}\rho\omega^2r^2$（轴心压强为零、边缘最大，用于分析离心铸造）；
上盖\textbf{边缘}开孔时 $p=\frac{1}{2}\rho\omega^2(r^2-R^2)$，
上盖内表面液体处于\textbf{真空}、中心压强最小（离心泵与风机靠此把流体从叶轮中心吸入）。

\textbf{常考题型：}选择（自由面形状、同一水平面压强分布方向）；
计算（赵琴书习题 2.7：求水正好与容器中心触底、顶部与容器同高时的角速度 $\omega$；
习题 2.8：分"上盖中心开孔"和"顶盖边缘开口"两种情况求压强表达式；
习题 2.9：充水深度 $h$ 的容器求不溢出水的极限转速）。

\textbf{重要度 \xstarfull{3}}\quad\yiju{E1 slide33「第二章＝计算题基础入门、重点掌握方法」；E4 历年真题中未见直接题号，仅章节级依据；E3「真题大部分源于课后习题改编」，赵琴书习题 2.7--2.9 全属本节}\quad\nandu{难}

\begin{banfa}
\textbf{①原点选择决定积分常数}，这是本节最大陷阱：
原点取在\emph{自由面中心}时 $C=p_0$、公式最简；
原点取在\emph{容器底面}时必须用"某处相对压强为零"重新定 $C$（有盖容器常用
"$r=0,z=H$ 处与大气接触"）。
\textbf{②体积守恒定 $\omega$}：旋转前后液体体积不变。
抛物面体积 $=\frac{1}{2}\pi R^2\Delta H$，用它列"原液面体积 $=$ 旋转后体积"的方程，
就能解 $\omega$ 或 $\Delta H$——赵琴书习题 2.7、2.9 都是这一个套路。\\
\textbf{③$a$ 与 $\omega$ 两类题共用一套流程：}写 $f_x,f_y,f_z$ $\to$ 积分求 $\dd p$
$\to$ 令 $\dd p=0$ 得等压面 $\to$ 体积守恒定位置 $\to$ 代求压强。
\emph{记流程比记公式更划算。}
\end{banfa}

\begin{yicuo}
\begin{enumerate}
\item 旋转容器中\textbf{同一水平面内压强不再相等}（随 $r$ 变化）。
  把绝对静止的"等深等压"结论无脑搬到旋转容器里，是本节最常见的失分点。\\
\item 自由面方程 $z_0=\dfrac{\omega^2r^2}{2g}$ 只有在\textbf{原点取在自由面中心}时成立；
  原点在别处时要改。\\
\item $\Delta H=\dfrac{\omega^2R^2}{2g}$ 是\textbf{最大}超高（在器壁 $r=R$ 处），
  不是平均超高；求某半径处的液面升高要把 $r$ 代进 $z_0$。
\end{enumerate}
\end{yicuo}
\end{kaodian}

\section{液体作用在平面上的总压力}

\begin{kaodian}{3}{平面图形的几何性质与常见图形的 $J_C$}
\textbf{是什么：}计算总压力与压力中心所需的纯几何量。

\begin{gongshi}{静矩与惯性矩的平行移轴定理（式 2.27、2.28）}
\[ \int_A y\,\dd A=y_C A,\qquad\qquad \int_A y^2\,\dd A=J_C+y_C^2A \]
\textbf{含义：}平面图形对 $x$ 轴的静矩等于\emph{面积}与\emph{形心到轴的距离}之积；
对 $x$ 轴的惯性矩等于\emph{对过形心平行轴的惯性矩} $J_C$ 加上 $y_C^2A$——
后一式即\textbf{平行移轴定理}，是推压力中心公式的关键。\\
\textbf{适用条件：}$y_C$ 必须从 $x$ 轴（一般取液面）量起；两轴必须平行。
\end{gongshi}

\begin{gongshi}{常见图形的形心与惯性矩（表 2.1）}
\begin{center}
\begin{tabularx}{\textwidth}{@{}lXXX@{}}
\toprule
图形 & 形心位置 & 面积 $A$ & 对过形心水平轴的 $J_C$ \\
\midrule
矩形（宽 $b$、高 $h$） & $h/2$ & $bh$ & $bh^3/12$ \\
三角形（底 $b$、高 $h$） & $2h/3$（自顶点量） & $bh/2$ & $bh^3/36$ \\
圆形（半径 $R$） & 圆心 & $\pi R^2$ & $\pi R^4/4$ \\
半圆（半径 $R$） & $4R/(3\pi)$ & $\pi R^2/2$ & $(9\pi^2-64)R^4/(72\pi)$ \\
圆环（外 $R$、内 $r$） & 圆心 & $\pi(R^2-r^2)$ & $\pi(R^4-r^4)/4$ \\
等边梯形（上 $a$、下 $b$、高 $h$） & $\frac{h(a+2b)}{3(a+b)}$ & $\frac{h(a+b)}{2}$ & $\frac{h^3(a^2+4ab+b^2)}{36(a+b)}$ \\
\bottomrule
\end{tabularx}
\end{center}
\textbf{含义／适用条件：}这些量\emph{只与图形几何有关}，与水深无关；
矩形、圆形的 $J_C$ 必须默写，梯形、半圆考场上可临时代入公式推导。
\end{gongshi}

\textbf{常考题型：}填空（默写矩形、圆形的 $J_C$）；计算（题目只给几何尺寸，
要求先求形心坐标再求总压力——表 2.1 就是"送分表"）。

\textbf{重要度 \xstarfull{3}}\quad\yiju{E1 slide33「第二章＝计算题基础入门、重点掌握方法」；E3「真题大部分源于课后习题改编」；赵琴书表 2.1 为 2.5、2.6 节全部计算的工具}\quad\nandu{易}
\end{kaodian}

\begin{kaodian}{5}{平面上的静水总压力与压力中心}
\textbf{是什么：}面积为 $A$ 的平面完全淹没在密度为 $\rho$ 的静止液体中，
平面上压强随深度线性变化，构成\textbf{非均匀分布力系}，可用一个合力代替。
由于各点压强都与平面正交，总压力 $P$ 也\textbf{垂直于平面}。

\begin{gongshi}{总压力大小——解析法（式 2.29）}
\[ P=\rho g\,y_C\sin\alpha\,A=\rho g h_C A=p_C A \]
\textbf{含义：}总压力等于\textbf{平面形心处压强}与其\textbf{淹没面积}之积。
注意 $h_C=y_C\sin\alpha$（$\alpha$ 为平面与液面夹角，$y_C$ 为形心沿平面的坐标），
$p_C$ 是形心压强，也等于被淹没面积上的\textbf{平均压强}。\\
\textbf{适用条件：}任意形状的平面；$p_C$ 必须用\textbf{相对压强}（即以大气压为零点），
否则要把大气压那部分单独处理掉。
\end{gongshi}

\begin{gongshi}{压力中心位置（式 2.30、2.31）}
\[ y_D=y_C+\frac{J_C}{y_C A} \]
\textbf{含义：}\textbf{压力中心 $D$} 是静水压力作用线与平面的交点，由对 $x$ 轴的力矩定理推出。
因为 $\dfrac{J_C}{y_C A}>0$ 恒成立，所以 $y_D>y_C$——
\textbf{压力中心永远位于平面形心之下}，两点沿平面的距离为 $J_C/(y_C A)$。\\
\textbf{适用条件：}平面\emph{倾斜或铅垂}（即不水平）；$y_C$、$y_D$ 都从\textbf{液面}量起。
若平面\textbf{水平放置}，压强处处相等，此时压力中心与形心\textbf{重合}（这是唯一特例）。
轴对称图形的压力中心一定在对称轴上，不必算 $x_D$。
\end{gongshi}

\begin{gongshi}{图解法（仅适用于一边平行于水面的矩形平面）}
\[ P=\Omega\cdot b \]
\textbf{含义：}总压力大小等于\textbf{压强分布图的面积} $\Omega$ 乘以受压平面的\textbf{宽度} $b$；
压力中心位置相当于压强分布图的\textbf{几何形心}。
\textbf{适用条件：}只能用于\emph{矩形}平面且其一边平行于水面；
绘制压强分布图时先标出受压面起点、终点的压强（沿水深线性），再以直线连接，
压强方向始终垂直于受压面。
\end{gongshi}

\textbf{例（赵琴书例 2.3）验证两种方法一致：}铅垂矩形闸门 $h=2\ \mathrm{m}$、$b=3\ \mathrm{m}$，
上缘距水面 $h_1=1\ \mathrm{m}$。图解法：梯形压强分布图面积 $\times b=P=117.6\ \mathrm{kN}$，
压力中心在梯形形心,距水面 $y_D=2.17\ \mathrm{m}$；解析法：
$P=\rho g\left(h_1+\dfrac{h}{2}\right)hb=9800\times2\times2\times3=117.6\ \mathrm{kN}$，
$y_D=y_C+\dfrac{J_C}{y_C A}=2+\dfrac{\frac{1}{12}\times3\times2^3}{2\times2\times3}=2.17\ \mathrm{m}$。
\emph{两法结果完全相同——考场上用图解法既快又能互相校验。}

\textbf{常考题型：}计算大题（铅垂或倾斜矩形闸门求总压力及作用点；
圆形闸门、三角形闸门求 $P$ 与 $y_D$）；证明题（赵琴书习题 2.12：证明圆形闸门水压力
对通过形心水平轴的力矩与形心水深 $h_C$ 无关）；填空（"压力中心永远
\underline{\ 低于\ }形心"）。

\textbf{重要度 \xstarfull{5}}\quad\yiju{E4 2015 年试题计算题含圆柱闸门静水总压力（需先求水平分力与作用点）；E1 slide33「第二章＝计算题基础入门、重点掌握方法」；E3「真题大部分源于课后习题改编」——赵琴书习题 2.10--2.13 全属本节；E2「第 1--3 章掌握基本公式与解释即可」}\quad\nandu{中}

\begin{banfa}
\textbf{①"两步走"标准解题步骤：}
(a) 定形心：画图找出受压面的形心位置，算出\textbf{形心淹深} $h_C$；
(b) 求大小：$P=\rho g h_C A$；
(c) 求作用点：$y_D=y_C+\dfrac{J_C}{y_C A}$，注意 $y_C$、$y_D$ 一律\textbf{从液面沿平面量}，
最后按题目要求换回"距某点的距离"。\\
\textbf{②图解法秒算梯形坡面：}铅垂矩形压强分布图必为\textbf{梯形}，
面积 $=\frac{1}{2}(\rho gh_1+\rho gh_2)\cdot h$，形心高度
$y_D=h_1+\dfrac{h(2\rho gh_2+\rho gh_1)}{3(\rho gh_1+\rho gh_2)}$——
熟记"上底、下底、高"的加权即可，比解积分快。\\
\textbf{③先判断平面是否水平}：水平面（如容器底）压力中心＝形心，
直接 $P=pA$ 就完事，不要再去套 $y_D=y_C+J_C/(y_C A)$。\\
\textbf{④闸门绕轴转动类题}（如扬压力、启门力）：先求 $P$ 与 $y_D$，
再对\emph{转轴}取力矩，把力矩平衡写出来解未知力——赵琴书习题 2.13 就是这一套路。
\end{banfa}

\begin{yicuo}
\begin{enumerate}
\item \textbf{$y_C$ 与 $h_C$ 分不清}：$h_C=y_C\sin\alpha$。
  平面铅垂（$\alpha=90^\circ$）时两者相等，倾斜时\emph{不相等}，这是倾斜平面题的头号错因。\\
\item \textbf{压力中心与形心混淆}：形心是\emph{几何}属性，压力中心是\emph{力学}（合力作用点）属性；
  除水平面外二者永不重合，且总有 $y_D>y_C$。\\
\item 图解法用在\emph{非矩形}或\emph{无平行边}的平面上——\textbf{错}，
  解法只对矩形（一边平行水面）成立。\\
\item 用绝对压强去算总压力会多出 $p_a A$ 这一项（除非题目明确要求考虑大气压），
  一般题默认不计大气压，用相对压强。
\end{enumerate}
\end{yicuo}
\end{kaodian}

\begin{kaodian}{4}{压强分布图（矩形、倾斜平面）}
\textbf{是什么：}在受压平面上按一定比尺绘制的、表示压强\textbf{大小与方向}分布的图形。
压强沿水深线性变化，故先标出受压面\textbf{起点和终点}的压强，再以\textbf{直线连接}，
压强方向\emph{始终垂直于受压面}。

\begin{gongshi}{三类典型受压面的压强分布图}
\begin{center}
\begin{tabularx}{\textwidth}{@{}lXXX@{}}
\toprule
受压面 & 分布图形状 & 关键结论 \\
\midrule
铅垂矩形（上缘在液面下） & \textbf{梯形} & 上边 $\rho gh_1$、下边 $\rho g(h_1+h)$，$\Omega$ 为梯形面积 \\[1mm]
铅垂矩形（上缘恰在液面） & \textbf{三角形} & 顶点在液面，底边 $\rho gh$，$\Omega=\frac{1}{2}\rho gh\cdot h$ \\[1mm]
水平面（容器底） & \textbf{矩形（等值）} & 压强处处相等，压力中心＝形心 \\[1mm]
倾斜平面 & 上小下大的\textbf{梯形}（或三角形） & 分布图仍画在受压面上，方向垂直受压面，长度表示压强大小 \\
\bottomrule
\end{tabularx}
\end{center}
\textbf{含义：}分布图只是"把压强沿水深线性变化的规律画出来"，
它既是图解法的输入，也是判断合力方向与作用点高低的直观工具。\\
\textbf{适用条件：}液体静止、自由面为水平面；压强用相对压强（液面处为零）。
\end{gongshi}

\textbf{常考题型：}作图题（画出闸门、坝面、容器壁上的压强分布图，
标出关键点数值与方向）；选择（"压强分布图呈怎样的形状"）；
计算（用图解法求总压力与作用点——赵琴书例 2.3、习题 2.11）。

\textbf{重要度 \xstarfull{4}}\quad\yiju{E1 slide33「第二章＝计算题基础入门、重点掌握方法」；E3「真题大部分源于课后习题改编」——赵琴书习题 2.11 明确要求"分别用解析公式和图解法"求解；E1 slide16「选择/判断考概念辨析」}\quad\nandu{易}

\begin{banfa}
\textbf{画图四句话：}① 先在受压面上定出\emph{上端}与\emph{下端}两点；
② 上端压强 $=\rho g\times$（上端淹深），下端同理；③ 两点\textbf{直线相连}（绝对静止时永远是直线）；
④ 箭头方向\textbf{垂直于受压面、指向受压面}。\\
\textbf{用图解法的合法性判断：}受\emph{气压}影响的、或受压面不是矩形的，不要用图解法。
遇到"坝面上有上游水又有下游水"的双面受压，
\emph{分别}画两侧分布图，最后把总压力相减（矢量和），别硬套一个梯形。
\end{banfa}
\end{kaodian}

\section{液体作用在曲面上的总压力}

\begin{kaodian}{5}{曲面总压力：水平分力与铅垂分力}
\textbf{是什么：}工程上最常见的曲面是具有水平或铅垂主轴的\textbf{圆柱形曲面}（二向曲面）。
曲面上不同水深处的压力\emph{方向不同}，\textbf{不能直接在曲面上积分求合力}，
必须先把微元压力 $\dd P=\rho gh\,\dd A$ 分解为水平、铅垂两个分量，再分别积分。

\begin{gongshi}{水平分力（式 2.32）}
\[ P_x=\int_A \rho gh\,\dd A\cos\alpha=\rho g h_C A_x \]
\textbf{含义：}\textbf{水平分力等于该曲面在铅垂投影面上的投影面 $A_x$ 所受到的总压力}。
$A_x$ 是曲面在 $yOz$ 面上的投影，$h_C$ 是其\textbf{形心的淹深}，
$A_x$ 上的压强分布为铅直线性，所以水平分力可用 2.5 节的平面总压力公式计算，
其作用线也\textbf{通过投影面 $A_x$ 的压力中心}。\\
\textbf{适用条件：}液体静止；投影面取\emph{垂直于所求水平分力}方向的平面；
对上下都有液体的曲面，要按两侧分别积分再合成。
\end{gongshi}

\begin{gongshi}{铅垂分力与压力体（式 2.33--2.35）}
\[ P_z=\rho g V,\qquad V=\int_A h\,\dd A_z,\qquad
   P=\sqrt{P_x^2+P_z^2},\qquad \alpha=\arctan\frac{P_z}{P_x} \]
\textbf{含义：}\textbf{铅垂分力等于压力体的液重}。$V$ 是以曲面 $ab$ 为底面、
以曲面在自由液面（或其延长面）上的投影面 $A_z$ 为顶面、
以及曲面周边各点向上投影的\textbf{所有铅垂母线}所围成的空间体积；
$P_z$ 的作用线\textbf{通过压力体的重心}。
总压力作用线必过 $P_x$、$P_z$ 两作用线的交点，但\textbf{这个交点不一定在曲面上}。\\
\textbf{适用条件：}\emph{压力体只是一个空间几何体积}，其内部是否有液体\emph{不影响}计算。
\end{gongshi}

\textbf{常考题型：}计算大题（弧形／圆柱闸门求 $P_x$、$P_z$、总压力与倾角，
再对转轴取矩求启闭力——赵琴书例 2.4、习题 2.14--2.17）；
填空（"曲面总压力的水平分力等于\underline{\ 曲面在铅垂投影面上的总压力\ }"；
"铅垂分力等于\underline{\ 压力体体积的液重\ }"）。

\textbf{重要度 \xstarfull{5}}\quad\yiju{E4 2015 年试题计算题含圆柱闸门静水总压力；E5 题库计算题第 1 题＝弧形闸门求 $F_x$、$F_z$ 与力矩；E1 slide33「第二章＝计算题基础入门、重点掌握方法」；E3「真题大部分源于课后习题改编」——赵琴书习题 2.14--2.17}\quad\nandu{难}

\begin{banfa}
\textbf{①固定"四步走"解题框架：}
(a) 画图，把曲面在\emph{铅垂面}上的投影画出来 $\to$ 算 $A_x$、$h_C$ $\to$ $P_x=\rho g h_C A_x$；
(b) 画\emph{压力体}：曲面 $\to$ 铅垂母线闭合 $\to$ 顶面取水面（或水面的延长面）$\to$ 算 $V$；
(c) 判方向：$P_z=\rho gV$，实压力体向下、虚压力体向上；
(d) 合成 $P=\sqrt{P_x^2+P_z^2}$、$\alpha=\arctan(P_z/P_x)$。
需要求力矩时\emph{再加一步}：对转轴分别取 $P_x$、$P_z$ 的力臂，写力矩平衡。\\
\textbf{②"求总压力"还是"求总压力与方向"要看清}：只写数值不给方向，大题会丢分。\\
\textbf{③圆弧面（圆柱面）的捷径：}圆弧面上每点的压强都沿半径指向圆心，
\emph{所有微元压力的作用线都通过圆心}——因此圆弧面总压力的作用线\textbf{必过圆心}。
赵琴书例 2.4 就用这一条直接断定合力过 $O$ 点，省掉求作用点的麻烦。\\
\textbf{④双侧受压（上下游都有水）}：分别算上游、下游的 $P_x$，
取差值；$P_z$ 则把两个压力体按虚实相抵的原则合成，再乘 $\rho g$。
赵琴书习题 2.15（圆柱闸门上下游水深 $H_1$、$H_2$ 不同）就是这一类。
\end{banfa}
\end{kaodian}

\begin{kaodian}{5}{压力体：虚体／实体与方向判定}
\textbf{是什么：}压力体是求曲面铅垂分力时构造出来的那个\textbf{空间几何体积}，
由受力曲面、液体自由表面（或其延长面）以及两者之间的铅垂面围成。
赵琴书口径：它是积分 $\int_A h\,\dd A_z$ 所得的一个体积，
是一个\textbf{纯数学的概念，与体积内有无液体无关}。

\begin{center}
\begin{tabularx}{\textwidth}{@{}lXXX@{}}
\toprule
类型 & 定义 & 符号 & 铅垂分力方向 \\
\midrule
\textbf{实压力体} & 压力体与形成压力的液体在曲面的\textbf{同侧} & $(+)$ & 铅垂\textbf{向下} \\
\textbf{虚压力体} & 压力体与形成压力的液体在曲面的\textbf{异侧} & $(-)$ & 铅垂\textbf{向上} \\
\bottomrule
\end{tabularx}
\end{center}

\textbf{压力体画法（四步，赵琴书未列此表但杨树人整理稿给出）：}
(a) 把受力曲面按情况\textbf{分成若干段}；
(b) 找出各段的\textbf{等效自由液面}；
(c) 画出每段的压力体并\textbf{判定虚实}；
(d) 按\textbf{虚实相抵}的原则把各段压力体合成，得到最终压力体。

\textbf{常考题型：}选择（"压力体的虚实与内部有无液体（有关／无关）"$\to$ \textbf{无关}）；
简答（压力体的定义、画法与虚实判定）；计算（先画压力体再算 $P_z$，
是曲面总压力题的中段核心）。

\textbf{重要度 \xstarfull{5}}\quad\yiju{E5 题库计算题第 1 题＝弧形闸门求 $F_x$、$F_z$ 与力矩（必须先正确取压力体）；E4 2015 年试题计算题含圆柱闸门静水总压力；E1 slide33「第二章＝计算题基础入门、重点掌握方法」}\quad\nandu{难}

\begin{banfa}
\textbf{①虚实一句话判断：}"\textbf{液体在里、体积在外}$\to$虚；液体在外、体积在里$\to$实"
更稳的说法是——\emph{把水从曲面处"填"满压力体，
如果填进去的水被曲面托在上面（水在曲面上方），则实压力体、力向下；
如果压力体是"抽出来的空腔"（水在曲面下方、压力体是曲面凸向上方空着的那块），
则虚压力体、力向上}。\\
\textbf{②判完方向再代大小}：$P_z=\rho gV$ 恒取正数，
方向由虚实单独判定，\emph{不要用负号去表达方向}，避免最后合成时符号弄乱。\\
\textbf{③压力体的"顶"一定取在自由液面（或其延长面）}，
不是取在曲面最高点，也不是取在容器壁。液面有多个/有压密闭容器时，
先用 $p=p_0+\rho gh$ 找到\textbf{等效自由液面}（即把 $p_0$ 折合成一个等效液柱高度），再画压力体。
\end{banfa}

\begin{yicuo}
\begin{enumerate}
\item "\textbf{压力体里必须真的充满液体}"——\textbf{错}。压力体是纯几何体积，
  赵琴书原文明确"与体积内有无液体无关"，虚压力体本身就是空的。\\
\item "\textbf{压力体越大铅垂分力越大}"——只有当虚实相同时才成立；
  有虚有实时要先做\textbf{矢量（虚实）合成}再乘 $\rho g$。\\
\item 圆弧闸门题里把 $\rho gV$ 的方向判反：只要记住
  "\emph{虚压力体、铅垂分力向上}"，赵琴书例 2.4 中压力体 $ABC$ 为虚压力体，
  故 $P_z$ 向上，与"闸门被水往上托"的直觉一致。\\
\item 压力体的顶面取错（取成曲面端点连线或容器壁），
  导致 $V$ 算成"半个弓形"而不是"弓形 $+$ 矩形"。
\end{enumerate}
\end{yicuo}
\end{kaodian}

\begin{kaodian}{5}{浮力与阿基米德原理}
\textbf{是什么：}潜体或浮体表面是\textbf{封闭曲面}，其上液体压力的\textbf{水平分力互相抵消}，
只剩铅垂分力，即浮力。

\begin{gongshi}{浮力与阿基米德原理}
\[ P_z=\rho g V_{\text{排}}=\rho g V_{\text{物浸}} \]
\textbf{含义：}作用在潜体（浮体）上的静水总压力只有铅垂方向的分量，方向\textbf{铅直向上}，
大小等于物体\textbf{所排开的同体积液体所受的重力}——这就是阿基米德原理。
浮力的作用线通过\textbf{排开液体的形心}，该点称为\textbf{浮心}。
（推导要点：物体上、下半表面在同一水平投影位置的压力方向相反、大小差为 $\rho gV$。）\\
\textbf{适用条件：}物体全部或部分浸没在静止液体中；液体均质不可压缩；
浮力只取决于物体\textbf{浸入部分的体积}，与物体材料和形状无关。
\end{gongshi}

\textbf{与压力体的联系：}浮体平衡时 $P_z$ 等于排开液体的重量；
对\emph{全部}浸没的物体，压力体因上下两半恰互相抵消而给出纯向上的合力。
工程应用：船舶浮性、浮力式液位计、潜浮与稳性计算。

\textbf{常考题型：}选择／判断（"浮力大小只与（物体体积／排开液体体积）有关"$\to$ 排开液体体积；
"浮力作用线通过（物体形心／排开液体形心）"$\to$ 排开液体形心即浮心）；
简答（用压力体解释浮力的产生）；计算（赵琴书例 2.5：圆柱形压力罐
求端盖、上下半圆筒所受水压力，并求连接螺栓所受总拉力 $T=P_{\text{上}}=49.54\ \mathrm{kN}$）。

\textbf{重要度 \xstarfull{5}}\quad\yiju{E1 slide33「第二章＝计算题基础入门、重点掌握方法」；E3「真题大部分源于课后习题改编」；E2「第 1--3 章掌握基本公式与解释即可」；浮力即曲面铅垂分力的封闭曲面对偶情形（赵琴书 2.6 节、例 2.5）}\quad\nandu{中}

\begin{banfa}
\textbf{①"浮力三问"速答模板：}多大 $\to$ $\rho gV_{\text{排}}$；向哪 $\to$ 铅直向上；
过哪点 $\to$ 浮心（排开液体的形心，\emph{不是}物体形心）。\\
\textbf{②螺栓／法兰受力题统一套路：}先算曲面（半球、封头）所受的铅垂分力 $P_z=\rho gV$，
由于对称，\emph{螺栓拉力恰好承担这一分量}——赵琴书例 2.5 中 $T=P_{\text{上}}=49.54\ \mathrm{kN}$。
注意这类题的内压由\textbf{压力表读数}给出，$p_m=\rho g h$ 换算成液柱高后再并入压力体。\\
\textbf{③有压密闭容器画压力体时}：先用 $h=\dfrac{p_m}{\rho g}$ 把表压折成\textbf{等效液柱高度}，
把自由液面"抬高"到虚构位置，再按普通敞口容器画压力体——这一步是赵琴书例 2.5 的关键技巧。
\end{banfa}

\begin{yicuo}
\begin{enumerate}
\item 认为"浮力作用线通过物体形心"——\textbf{错}，通过\emph{排开液体的形心}（浮心）；
  两者只有在物体均质且全浸没时才重合。\\
\item 把浮力当作"物体自重的一部分"——\textbf{错}，浮力是\emph{液体对物体表面压力}的合力，
  与物体自重是两个不同的力，两者相等只在漂浮／悬浮平衡时成立。\\
\item 认为"液体深度越大浮力越大"——\textbf{错}，浮力只与排开液体体积有关，
  完全浸没后与深度无关（这正是阿基米德原理的核心）。
\end{enumerate}
\end{yicuo}
\end{kaodian}

\section{本章小结}

\begin{center}
\begin{tabularx}{\textwidth}{@{}lcXX@{}}
\toprule
考点 & 重要度 & 难度 & 一句话抓住 \\
\midrule
绝对静止与相对静止 & \xstarfull{4} & 易 & 都由欧拉方程管；只差质量力里有无惯性力 \\
静压强特性与单位 & \xstarfull{4} & 易 & 内法线＋与方位无关；$1\ \mathrm{bar}=10^5\ \mathrm{Pa}$ \\
欧拉平衡微分方程 & \xstarfull{4} & 中 & $\grad p=\rho\bF$；$\dd p=\rho(f_x\dd x+f_y\dd y+f_z\dd z)$ \\
等压面及性质 & \xstarfull{4} & 易 & 等压面$\perp$质量力；$f_x\dd x+f_y\dd y+f_z\dd z=0$ \\
静力学基本方程与水头 & \xstarfull{5} & 中 & $z+\dfrac{p}{\rho g}=C$，测压管水头线水平 \\
压强表示法 & \xstarfull{4} & 易 & $p_v=p_a-p_{\text{abs}}=-p_M$，虹吸管顶是真空 \\
静压强的测量 & \xstarfull{4} & 中 & 找最低\emph{同种}液体等压面，向下加向上减 \\
等加速直线运动容器 & \xstarfull{3} & 中 & $\tan\theta=a/g$，等压面为斜面，靠体积守恒定位 \\
等角速度旋转容器 & \xstarfull{3} & 难 & $z_0=\dfrac{\omega^2r^2}{2g}$，$\Delta H=\dfrac{\omega^2R^2}{2g}$，水平面内不等压 \\
平面几何性质与 $J_C$ & \xstarfull{3} & 易 & 矩形 $bh^3/12$、圆 $\pi R^4/4$ \\
平面总压力与压力中心 & \xstarfull{5} & 中 & $P=\rho gh_C A$；$y_D=y_C+\dfrac{J_C}{y_C A}$（恒低于形心） \\
压强分布图 & \xstarfull{4} & 易 & 铅垂矩形＝梯形／三角形；图解法只对矩形 \\
曲面总压力（$P_x$、$P_z$） & \xstarfull{5} & 难 & $P_x=\rho gh_C A_x$；$P_z=\rho gV$ \\
压力体虚实与方向 & \xstarfull{5} & 难 & 纯几何体积，与有无液体无关；异侧为虚、力向上 \\
浮力与阿基米德原理 & \xstarfull{5} & 中 & $P_z=\rho gV_{\text{排}}$，向上，过浮心 \\
\bottomrule
\end{tabularx}
\end{center}
